\documentclass[10pt,letterpaper]{article}
\usepackage[hmargin=1in,vmargin=0.85in]{geometry}
\date{}
\usepackage{xurl}
\setlength{\emergencystretch}{2em}
\title{Distinguishing a Missing Freeway Detector Reading from Emerging Congestion}
\author{Arjun Vijay Prakash \\ \small ITE Annual Meeting 2027 | Transportation Operations |
Abstract draft}

\begin{document}

\maketitle
\begin{abstract}
When a freeway detector stops reporting, an operator has to decide whether to inspect the
equipment or look for a developing queue. CREST presents the last valid speed and its age
alongside neighboring detectors, a persistence baseline, and a graph-based reconstruction. The
tool suggests a detector check when neighboring stations remain clear, a queue check when
adjacent stations report concurrent low speeds, and an unavailable status when the evidence is
insufficient.
\end{abstract}

\section*{Case demonstration}
On southbound I-880, we hid a five-minute Caltrans PeMS speed reading at station 401362. The
tool suggested inspecting the detector, and its last valid speed was one interval old. The
withheld reading gave an absolute error of 1.77 km/h for persistence and 1.85 km/h for the graph
reconstruction. In a separate window with concurrent low speeds, the tool suggested
corroborating a queue. The District 7 replay also suggested inspecting the detector for an
injected outage on westbound I-10.

\section*{What the replay shows}
In the expanded archive replay, graph reconstruction reduced speed MAE from 5.86 to
4.65 km/h on 468 scored target instances, even though persistence was slightly better
in the selected I-880 case. Both methods missed all ten low-speed queue-proxy positives.
The demonstration lets an operator trace a suggestion to the readings available at the
time and score an estimate against a withheld value.
\end{document}
